% part: intuitionistic-logic
% chapter: propositions-as-types
% section: terms-to-proofs

\documentclass[../../../include/open-logic-section]{subfiles}

\begin{document}

\olfileid[id]{int}{pty}{tp}

\olsection{Memulihkan \usetoken{P}{derivation} dari Term Bukti}

Sekarang mari kita tinjau arah sebaliknya: menerjemahkan term bukti~$N$
kembali menjadi !!a{proof} dalam \Log{tN2i}. Secara umum, hal ini hanya dapat
dilakukan relatif terhadap konteks yang memuat pasangan $x:!A$ bagi setiap
variabel~$x$ yang muncul bebas dalam term bukti tersebut. Namun, kita mulai
dengan contoh tanpa variabel bebas, yaitu term
\[
\lambd[\typeof{x}{!A \land
  !B}][\pair{\proj{2}{x}}{\proj{1}{x}}]
\]
Term ini berbentuk $\lambd[\typeof{x}{!A \land !B}][N_1]$. Karena satu-satunya
aturan \Log{tN2i} yang menghasilkan term berbentuk demikian adalah
\Intro{\lif}, kita mengetahui bahwa !!{proof} yang dikodekan oleh~$N$ harus
berakhir dengan~\Intro{\lif}, yakni berakhir sebagai
  \[
  \Axiom$x : !A\land !B \fCenter \pair{\proj{2}{x}}{\proj{1}{x}} :
    !C$
  \RightLabel{\Intro\lif}
  \UnaryInf$\fCenter \lambd[\typeof{x}{!A \land
  !B}][\pair{\proj{2}{x}}{\proj{1}{x}}] : (!A \land !B) \lif !C$
  \DisplayProof
  \]
Kita belum mengetahui apa itu $!C$, tetapi kita mengetahui bahwa antesedennya
adalah $!A \land !B$ dan bahwa konteks premis harus memuat
$x:!A \land !B$.

Term bukti yang terkait dengan premis kini sudah ditentukan, yaitu
$\pair{\proj{2}{x}}{\proj{1}{x}}$. !!^{proof} yang dikodekannya harus
berakhir dengan \Intro{\land}. Kita masih belum mengetahui apa itu $!C$,
tetapi karena $!C$ merupakan kesimpulan \Intro{\land}, kita mengetahui bahwa
$!C$ harus berupa konjungsi, yakni $!C \ident (!C_1 \land !C_2)$. Kita
melanjutkan rekonstruksinya:
  \[
  \Axiom$x : !A\land !B \fCenter \proj{2}{x} : !C_1$
  \Axiom$x : !A\land !B \fCenter \proj{1}{x} : !C_2$
  \RightLabel{\Intro\land}
  \BinaryInf$x : !A\land !B \fCenter \pair{\proj{2}{x}}{\proj{1}{x}} :
    !C_1 \land !C_2$
  \RightLabel{\Intro\lif}
  \UnaryInf$\fCenter \lambd[\typeof{x}{!A \land
  !B}][\pair{\proj{2}{x}}{\proj{1}{x}}] : (!A \land !B) \lif (!C_1 \land !C_2)$
  \DisplayProof
  \]
Term bukti $\proj{2}{x}$ berarti bahwa sekuen teratas di sebelah kiri
diperoleh dengan \Elim{\land}[2], yakni premisnya berbentuk
$!D \land !C_1$. Di sebelah kanan, kita dapat menentukan bahwa inferensi
yang menghasilkan $!C_2$ harus berupa \Elim{\land}[1], dan premisnya
berbentuk $!C_2 \land !E$.
  \[
  \Axiom$x : !A\land !B \fCenter x: !D\land !C_1$
  \RightLabel{\Elim\land[2]}
  \UnaryInf$x : !A\land !B \fCenter \proj{2}{x} : !C_1$
  \Axiom$x : !A\land !B \fCenter x: !C_2\land !E$
  \RightLabel{\Elim\land[1]}
  \UnaryInf$x : !A\land !B \fCenter \proj{1}{x} : !C_2$
  \RightLabel{\Intro\land}
  \BinaryInf$x : !A\land !B \fCenter \pair{\proj{2}{x}}{\proj{1}{x}} :
    !C_1 \land !C_2$
  \RightLabel{\Intro\lif}
  \UnaryInf$\fCenter \lambd[\typeof{x}{!A \land
  !B}][\pair{\proj{2}{x}}{\proj{1}{x}}] : (!A \land !B) \lif (!C_1 \land !C_2)$
  \DisplayProof
  \]
Karena term bukti dari sekuen-sekuen teratas kini berupa variabel, kita
mengetahui bahwa sekuen-sekuen itu harus merupakan aksioma. Hal ini
menentukan $!C_1$, $!C_2$, $!D$, dan $!E$: harus berlaku
$!D \ident !A$ dan $!C_1 \ident !B$, sebab jika tidak, sekuen kiri bukan
aksioma; demikian pula harus berlaku $!C_2 \ident !A$ dan $!E \ident !B$,
sebab jika tidak, sekuen kanan bukan aksioma. Kita telah memulihkan bukti
secara lengkap:
  \[
  \Axiom$x : !A\land !B \fCenter x: !A\land !B$
  \RightLabel{\Elim\land[2]}
  \UnaryInf$x : !A\land !B \fCenter \proj{2}{x} : !B$
  \Axiom$x : !A\land !B \fCenter x: !A\land !B$
  \RightLabel{\Elim\land[1]}
  \UnaryInf$x : !A\land !B \fCenter \proj{1}{x} : !A$
  \RightLabel{\Intro\land}
  \BinaryInf$x : !A\land !B \fCenter \pair{\proj{2}{x}}{\proj{1}{x}} :
    !B \land !A$
  \RightLabel{\Intro\lif}
  \UnaryInf$\fCenter \lambd[\typeof{x}{!A \land
  !B}][\pair{\proj{2}{x}}{\proj{1}{x}}] : (!A \land !B) \lif (!B \land !A)$
  \DisplayProof
  \]

\end{document}
